\section{Introduction}\label{sec:introduction}

The smooth realization problem asks which measure-preserving systems can
be represented by smooth dynamics on a compact manifold.
Here a \emph{smooth positive-volume model} of an invertible
probability-preserving transformation $(X,\mu,T)$ consists of a compact
smooth manifold $M$ of finite dimension, a smooth positive probability
volume $\nu$, and a $C^\infty$ diffeomorphism $S$ preserving $\nu$,
together with a measure-preserving isomorphism
\[
 F:X_0\longrightarrow M_0,\qquad F\circ T=S\circ F,
\]
on invariant conull measurable sets. Both $F$ and $F^{-1}$ are measurable;
no continuity is required. Here smooth positive volume means a probability
measure with strictly positive smooth local densities,
$d\nu=f(x)|dx_1\cdots dx_d|$, where $f>0$ is smooth.
The manifold may have a smooth boundary, with these densities smooth up
to the boundary. No orientation of $M$ or orientation preservation by
$S$ is required. The manifold and its dimension may depend on the
original system.

A standard nonatomic probability space means a probability space
isomorphic, modulo null sets, to the unit interval with Lebesgue
probability.

Finite entropy is a necessary restriction in the smooth realization
problem. For example, Ruelle's differentiable entropy inequality
\cite[Theorem~2]{Ruelle1978} gives finite entropy from compact derivative
bounds. The question addressed here is
whether this entropy restriction suffices when the invariant measure
must be a smooth positive density.

For invertible ergodic transformations of standard nonatomic
probability spaces, two positive realization results help locate this
restriction.
Lind and Thouvenot realized finite-entropy ergodic transformations by
Lebesgue-preserving homeomorphisms of the two-torus
\cite{LindThouvenot1977}. They also proved that a hyperbolic
two-torus automorphism represents every such system with measure entropy
below the automorphism's topological entropy, using a suitable invariant
Borel probability of full support; Quas and Soo recall this earlier result in
\cite[\S1]{QuasSoo2016}. Quas and Soo's Theorem~2 extends the result to
toral automorphisms with no root-of-unity eigenvalue.
Choosing a sufficiently high iterate of a hyperbolic automorphism
therefore gives a smooth model for any prescribed finite-entropy
system, with an invariant Borel probability of full support.
Full support does not require a smooth positive density.

The requirement on the invariant measure is essential. A realization
with an arbitrary invariant Borel probability, which may be singular,
does not by itself give a smooth positive-volume model. This distinction is part of the
classical formulation discussed by Katok and Thouvenot
\cite[pp.~324--326]{KatokThouvenot1997}. They developed slow entropy
invariants and constructed ergodic $\Z^2$ actions having no Lipschitz
realization on a compact space of finite box dimension, even when the
model may carry an arbitrary invariant Borel probability
\cite[Proposition~4 and Corollary~2]{KatokThouvenot1997}.
The commuting generators must be realized together in that problem;
its obstruction does not by itself answer the question for one
transformation. Foreman and Weiss
\cite[p.~2607]{ForemanWeiss2022} also discuss finite-entropy realization,
in a formulation allowing an invariant measure merely equivalent to
volume. Our conclusion concerns the smooth positive-density version;
we do not assert an obstruction for every equivalent-to-volume
invariant measure. Their anti-classification theorems concern the complexity of the
measure-isomorphism relation among smooth transformations, including
an ergodic area-preserving torus subclass
\cite[Theorems~1--2 and Corollary~4]{ForemanWeiss2022}.
They do not assert that a particular abstract finite-entropy system
lacks such a model. Subsequent anti-classification results treat
weakly mixing smooth transformations \cite[Theorem~B]{Kunde2024} and mixing
zero-entropy smooth transformations \cite[Theorem~2]{GerberKunde2025Mixing};
these too study equivalence inside families of smooth models.

We give a negative answer to the smooth-volume realization question,
already among systems of zero entropy.
For a finite probability vector $p$, its Shannon entropy is
\[
 H(p)=-\sum_a p(a)\log_2 p(a),\qquad 0\log_2 0=0.
\]
The entropy $H_\mu(\mathcal P)$ of a finite measurable partition is
the entropy of its atom probabilities; the entropy of a finite random
variable is likewise the entropy of its law.
The Kolmogorov--Sinai entropy of $T$ is
\[
 h_\mu(T)=\sup_{\mathcal P}\lim_{n\to\infty}
 \frac1n H_\mu\!\left(\bigvee_{j=0}^{n-1}T^{-j}\mathcal P\right),
\]
where $\mathcal P$ ranges over finite measurable partitions.

\begin{theorem}\label{thm:main}
There exists an ergodic invertible measure-preserving transformation
$(X,\mu,T)$ of a standard nonatomic probability space such that
$h_\mu(T)=0$ and $(X,\mu,T)$ is not measurably conjugate, on invariant conull sets, to
any $C^\infty$ diffeomorphism of a compact finite-dimensional manifold
preserving a strictly positive smooth probability density. The manifold
may have smooth boundary and need not be orientable.
\end{theorem}

The proof isolates a property of smooth positive volume that remains
useful after an arbitrary measurable change of coordinates.
At a dyadic time $n=2^k$, partition the manifold into small cells,
choose a cell with its volume as probability, and choose two independent
points $x,y$ from its normalized volume. On most such pairs, one can
find a point $z$ whose future shadows that of $x$ and whose past shadows
that of $y$. More than the existence of this point is needed:
the law of $z$ on the retained pairs is bounded by twice the original
volume. The number of cells is at most $\exp(CB_kn)$ for a sequence
with $\sum_k B_k^{16}<\infty$.
The precise statement is \Cref{thm:switching}.

This switch estimate is the main analytic ingredient. Its proof begins
with a square-summable dyadic estimate for stationary diagonal
derivative logs, rather than an assumed rate in an ergodic theorem.
Finite-interval triangular changes of basis give the linear
estimates needed to solve a boundary problem joining the two orbit
halves. A relative exterior-power estimate controls the change in
small endpoint minors, and a determinant identity and global
injectivity control the law of the switch.
The quantitative summability and the bound on the output measure are
the features that will obstruct realization.

We construct the opposing symbolic system from finite sets
$\mathcal W_i$ of words of lengths $m_i$. Each word at the next level is
a concatenation of words at the preceding level, and the system
records their compatible alignment phases. Word labels are nonzero
vectors over $\F_2$. Certain pairs of labels separated by a power-of-two
number of blocks must be orthogonal for a fixed alternating form.
These constraints coexist with high entropy of short label lists and
with a positive symbol-distance separation between distinct or
misaligned words.

The entropy of the label lists makes two independent lists unlikely
to satisfy most of the orthogonality tests. A smooth switch would
force their phases and most of their labels to agree with the
corresponding past and future lists of a single legal symbolic point.
Those lists do satisfy the tests. Summability of the smooth partition
cost, together with the growth of the word scales, supplies times at
which conditioning on a cell loses too little entropy to reconcile
the two conclusions.

The order of choices is important: the symbolic system is fixed first.
Only then is a hypothetical smooth model, with any one finite
dimension, considered. The scale argument absorbs all constants
depending on that model. Thus the construction does not select a
different counterexample for each dimension.

\paragraph{Organization.}
\Cref{sec:stationary,sec:linear,sec:switching} prove the switching
estimate. \Cref{sec:words} constructs the finite word families.
\Cref{sec:process} supplies an ergodic probability on their phase
extension and establishes the window entropy and graph constraints.
\Cref{sec:obstruction} compares the two laws of one graph-passing
statistic and proves \Cref{thm:main}.

\paragraph{Conventions.}
Logarithms in entropies and combinatorial estimates are to base $2$;
natural logarithms are written $\ln$, and derivative growth is written
with $\exp$ or $e$. Total variation distance between finite laws $p,q$
is $\TV(p,q)=\tfrac12\sum_a|p(a)-q(a)|$.
For finite random variables, conditional entropy means
\[
 H(Y\mid Z)=\sum_{z:\Pr(Z=z)>0}\Pr(Z=z)\,
 H\bigl(\mathcal L(Y\mid Z=z)\bigr).
\]
Conditional partition entropy uses the same definition for atom labels.
We write $H_{\mathrm{bin}}(s)=-s\log_2 s-(1-s)\log_2(1-s)$
for binary entropy, with the same zero-mass convention at the endpoints.
Constants denoted by $C$ may change from line to line.
In the smooth part they depend only on the fixed smooth data,
dimension, and a fixed exceptional-mass tolerance, unless otherwise
specified. All asymptotic estimates there are uniform over the
points, bins, and finite-interval frames satisfying the displayed
hypotheses. The empty determinant and empty product equal one.
